\documentclass[11pt]{amsart}
\usepackage[margin=1.1in]{geometry}
\usepackage{amssymb,amsmath,mathtools}
\usepackage{longtable,booktabs}
\usepackage[hidelinks]{hyperref}

\newtheorem{theorem}{Theorem}
\newtheorem{lemma}[theorem]{Lemma}
\newtheorem{proposition}[theorem]{Proposition}
\newtheorem{corollary}[theorem]{Corollary}
\theoremstyle{definition}
\newtheorem{definition}[theorem]{Definition}
\newtheorem{remark}[theorem]{Remark}
\newtheorem{question}[theorem]{Question}

\newcommand{\Z}{\mathbb{Z}}
\newcommand{\M}{M}
\newcommand{\kmax}{k_{\max}}

\title[Rainbow 4-cycles with $n$ colours]{Exact small values and cyclic constructions\\
for a rainbow 4-cycle problem of Burr, Erd\H{o}s, Graham and S\'os}
% TODO(author): author list, affiliations, and whether/how to disclose AI assistance.
\author{[Author]}
\date{Draft, \today}

\begin{document}

\begin{abstract}
Let $\M(n)$ be the largest number of edges of an $n$-vertex graph whose edges can be
coloured with $n$ colours so that every $4$-cycle receives four distinct colours.
Burr, Erd\H{o}s, Graham and S\'os asked whether $\M(n)\ge\varepsilon n^2$ for some
$\varepsilon>0$ (Erd\H{o}s problem \#810), and believed not. We determine $\M(n)$ for all
$n\le 12$, the new values being $\M(10)=30$, $\M(11)=34$ and $\M(12)=40$, by exhaustive
enumeration, independently implemented several times, together with SAT proofs that are
checked by a DRAT proof checker for $n\le 11$. This confirms that the lower bounds found by
AlphaEvolve for $n\le 12$ are optimal.
We give explicit certificates improving the known lower bounds for $14\le n\le 25$ and
extending them to $n\le 128$. Many of the best constructions for $26\le n\le 64$ are
\emph{cyclic orbit colourings}: Cayley graphs on $\Z_n$ whose colouring is equivariant under
rotation. For this family we compute the exact optimum for $12\le n\le 41$, $n\ne35$; it
equals exactly $n/4$ orbits, that is $n^2/4$ edges, for $n=12,16,\dots,40$. We also prove
that offset rules that are linear, or piecewise affine with a bounded number of pieces, give
only $o(n)$ orbits. The data are consistent with the expected negative answer.
\end{abstract}

\maketitle

% =====================================================================
\section{Introduction}

For graphs $G$ and $L$, an edge colouring of $G$ is \emph{$L$-rainbow} if every copy of
$L$ in $G$ (not necessarily induced) has all its edges coloured differently. Burr,
Erd\H{o}s, Graham and S\'os~\cite{BEGS89} introduced
\[
  \chi_S(n,e,L)=\min\{r:\ \text{some $n$-vertex graph with $e$ edges has an $L$-rainbow $r$-colouring}\}
\]
and studied it for complete graphs, paths, cycles and other small graphs. For $L=C_4$ the natural
threshold question is how many edges an $n$-vertex graph can have if $n$ colours suffice:
\[
  \M(n)=\max\{e:\ \chi_S(n,e,C_4)\le n\}.
\]

\begin{question}[{\cite[p.~273]{BEGS89}; Erd\H{o}s problem \#810~\cite{EP810}}]
Is there $\varepsilon>0$ such that $\M(n)\ge\varepsilon n^2$ for all large $n$?
\end{question}

Burr, Erd\H{o}s, Graham and S\'os believed that the answer is negative.
They noted that $\M(n)\ge c\,g(n;7,4)$, where $g(n;7,4)$ is the largest number of triples
on $n$ points with no $7$ points spanning $4$ triples. Whether $g(n;7,4)=o(n^2)$ is itself
open (Erd\H{o}s problem \#1178). For $P_4$ in place of $C_4$ they showed that the answer is
negative, via the Ruzsa--Szemer\'edi theorem~\cite{RS78}.
Sárközy and Selkow~\cite{SS06} proved that $\chi_S(n,\varepsilon n^2,L)/n\to\infty$ for every
connected bipartite $L$ that is neither a star nor complete bipartite. This excludes
$C_4=K_{2,2}$, and in~\cite{BEGS89} the case $C_4$ is called ``one of the most vexing questions''.
Tao~\cite{TaoForum} observed that a negative answer would imply a finite-field
Ajtai--Szemer\'edi theorem, via the colouring $(a,b)\mapsto a+b$ of a bipartite graph on
$\mathbb F_2^d\times\mathbb F_2^d$.

On the problem's discussion forum, Tao~\cite{TaoAE} reported dense constructions for
$4\le n\le 25$ found with AlphaEvolve~\cite{AE25}, with no guarantee of optimality. He
suggested computing exact values with SAT solvers. This note does that, and also pushes the
constructions further.

\subsection*{Results}
\begin{enumerate}
\item \textbf{Exact values} (Theorem~\ref{thm:exact}): $\M(n)$ for every $n\le12$
  (Table~\ref{tab:exact}), including $\M(10)=30$, $\M(11)=34$, $\M(12)=40$.
  We also give the number of extremal graphs for $7\le n\le 11$.
\item \textbf{Lower bounds} (Table~\ref{tab:lb}): explicit certificates for every
  $13\le n\le 64$ and many $n\le128$. They improve the AlphaEvolve values for all
  $14\le n\le 25$; for example $\M(21)\ge126=0.286\,n^2$.
\item \textbf{Cyclic orbit colourings} (Section~\ref{sec:cyclic}): we compute the exact
  optimum $\kmax(n)$ of this family for $12\le n\le 41$, $n\ne 35$.
  \begin{itemize}
  \item $\kmax(4m)=m$ for $3\le m\le 10$, so $\M(n)\ge n^2/4$ at $n=12,16,\dots,40$.
  \item The number of optimal solutions at $k=n/4$ collapses as $n$ grows, to $4$
    (up to symmetry) at $n=40$, and no $n/4$ solution has been found for $44\le n\le100$.
  \item Proposition~\ref{prop:linear}: rules with linear or boundedly piecewise-affine
    offsets give only $o(n)$ orbits.
  \end{itemize}
\end{enumerate}

All certificates, programs, solver logs and proof-checker outputs are in a public
repository~\cite{Repo} (see Appendix~\ref{app:repro}). Every certificate quoted here is
re-verified by an independent, deliberately simple checker when the tables are built.

% =====================================================================
\section{Preliminaries}

All graphs are simple. A colouring \emph{uses $k$ colours} if it takes values in
$\{0,\dots,k-1\}$. Call $G$ \emph{$k$-good} if it has a $C_4$-rainbow colouring with $k$
colours; thus $\M(n)$ is the largest number of edges of an $n$-good graph on $n$ vertices.

\begin{definition}
The \emph{conflict graph} $Q(G)$ has vertex set $E(G)$. Two distinct edges are adjacent in
$Q(G)$ if they lie on a common $4$-cycle of $G$. Concretely:
\begin{itemize}
\item adjacent edges $uv,vw$ conflict iff $u$ and $w$ have a common neighbour other than $v$;
\item disjoint edges $uv,xy$ conflict iff $ux,vy\in E(G)$ or $uy,vx\in E(G)$.
\end{itemize}
\end{definition}

A colouring is $C_4$-rainbow exactly when it is a proper colouring of $Q(G)$. Hence $G$ is
$k$-good iff $\chi(Q(G))\le k$.

\begin{lemma}\label{lem:hered}
If $G$ is $k$-good then so is every subgraph of $G$.
In particular, $\M(n)=e$ iff some $e$-edge $n$-vertex graph is $n$-good and no
$(e+1)$-edge $n$-vertex graph is $n$-good.
\end{lemma}

\begin{proof}
Every $4$-cycle of a subgraph is a $4$-cycle of $G$. The second statement follows because
deleting edges preserves goodness.
\end{proof}

Trivially $\M(n+1)\ge \M(n)+1$: add a pendant edge, which lies on no $4$-cycle.
Also $\M(n)\ge \mathrm{ex}(n,C_4)$, since a $C_4$-free graph can use a single colour.

% =====================================================================
\section{Exact values for $n\le12$}

\begin{theorem}\label{thm:exact}
The values of $\M(n)$ for $n\le 12$ are as in Table~\ref{tab:exact}.
\end{theorem}

\begin{table}[ht]
\centering
\caption{Exact values. ``Extremal'' is the number of $n$-good graphs with $\M(n)$ edges up
to isomorphism. E = exhaustive enumeration (\S\ref{sec:enum}), S = SAT
(\S\ref{sec:sat}), L = Lean kernel check~\cite{Repo}.}
\label{tab:exact}
\begin{tabular}{rrrp{0.55\textwidth}}
\toprule
$n$ & $\M(n)$ & extremal & upper bound $\M(n)+1$ excluded by\\
\midrule
0--1 & 0 & -- & trivial\\
2 & 1 & -- & trivial\\
3 & 3 & -- & trivial\\
4 & 5 & -- & L, S\\
5 & 7 & -- & L, S\\
6 & 11 & -- & L, E, S\\
7 & 14 & 5 & E, S\\
8 & 17 & 45 & E, S\\
9 & 23 & 3 & E, S (DRAT)\\
10 & 30 & 2 & E, S (DRAT, all cubes)\\
11 & 34 & 131 & E, S (DRAT, all cubes)\\
12 & 40 & -- & E (pruned, 2 checkers; unpruned, 4 checkers); S (DRAT, 9 of 12 cubes)\\
\bottomrule
\end{tabular}
\end{table}

The lower bounds are explicit certificates. For $n\le 12$ they were kernel-checked in Lean
with \texttt{decide}~\cite{Repo}, and the AlphaEvolve constructions~\cite{TaoAE} reach the
same values. The work is in the upper bounds, which we establish by three independent routes.

\subsection{Pruned orderly generation}\label{sec:enum}
We enumerate graphs up to isomorphism with \texttt{geng}~\cite{nauty}.
\begin{itemize}
\item \textbf{Pruning.} We use its \texttt{PRUNE} hook to reject every intermediate graph
  that is not $n$-good, where $n$ is the final number of vertices. \texttt{geng} adds vertices
  one at a time, so each intermediate graph is an induced subgraph of all its descendants. By
  Lemma~\ref{lem:hered} the pruning is therefore sound: no $n$-good output graph is lost.
\item \textbf{Checkers.} Goodness is decided by two independently written programs.
  Checker~A builds $Q(G)$ as bit masks, strips its $n$-core, and runs a clique bound followed
  by exact DSATUR. Checker~B enumerates the $4$-cycles explicitly and backtracks on the
  most-constrained edge.
\item \textbf{Validation.} Pruned and unpruned runs give identical sets of good graphs for
  $6\le n\le 11$, at $e=\M(n)$ and $e=\M(n)+1$. Total class counts produced by \texttt{geng} agree with a
  P\'olya count, for example $71{,}318$, $24{,}630{,}887$, $15{,}108{,}047$ and
  $2{,}411{,}961{,}516$ isomorphism classes for $(n,e)=(10,31),(11,34),(11,35),(12,41)$.
\item \textbf{Main run.} The run for $(12,41)$ used $512$ shards and $9{,}165$ CPU-seconds,
  and output no graph. Checker~B was run separately on the graphs of maximum degree at most
  $9$ (also all $512$ shards, no output). The cases $\Delta\in\{10,11\}$ are covered by the
  SAT certificates of \S\ref{sec:sat}.
\end{itemize}

\subsection{Unpruned enumeration with four independent checkers}
To remove any dependence on the pruning, the space of $(12,41)$ graphs was split by
\texttt{geng}'s residue classes into four parts.
\begin{itemize}
\item \textbf{Checkers.} Each part was processed by a separately written checker, with no
  shared code.
\item \textbf{Validation.} Each checker first reproduced the number of good graphs at
  $e=\M(n)$ and $e=\M(n)+1$ for $7\le n\le 11$: namely $5,45,3,2,131$ good graphs at
  $e=\M(n)$ and none at $e=\M(n)+1$.
\item \textbf{Totals.} The four parts contained $605{,}033{,}974$, $603{,}779{,}517$,
  $601{,}165{,}216$ and $601{,}982{,}809$ graphs. These sum to exactly the P\'olya count
  $2{,}411{,}961{,}516$.
\item \textbf{Outcome.} None of the graphs is $12$-good. All but $753$ of them contain a
  $13$-clique in $Q(G)$, found by exact branch and bound; the remaining $753$ were refuted by
  exhaustive colouring search.
\end{itemize}

\subsection{SAT with checked proofs}\label{sec:sat}
For given $n$ and $e$, the encoding uses three groups of variables.
\begin{itemize}
\item Edge variables $x_{uv}$ and colour variables $y_{uv,c}$, with at least one colour per
  present edge.
\item For each $4$-cycle $C$ and each pair $\{f,f'\}$ of its edges, a variable $z_{f,f'}$
  forced true when the other two edges of $C$ are present. Then
  $\neg z_{f,f'}\vee\neg y_{f,c}\vee\neg y_{f',c}$ for every colour $c$.
\item A totalizer for $\sum x_{uv}\ge e$.
\end{itemize}

Symmetry is broken by four constraints.
\begin{enumerate}
\item Vertex $0$ has maximum degree, and its neighbours come first.
\item Degrees are non-increasing within the neighbourhood of $0$ and within its complement.
\item A lexicographic row comparison for transpositions of equal-degree vertices in the same
  block~\cite{CMPS19}.
\item First-occurrence ordering of colours.
\end{enumerate}
Vertex relabelling is chosen from $G$ alone and colour relabelling afterwards, which makes
the combination sound. We also tested soundness exhaustively: every one of the $127$ classes
of $6$-good graphs on $6$ vertices has an accepted labelling, and no non-good graph is
accepted.

The search is split into \emph{cubes} by the maximum degree $\Delta$, from $0$ to $n-1$.
Each cube was solved by CaDiCaL~\cite{CaDiCaL}, and its DRAT proof was checked by
\texttt{drat-trim}~\cite{dratTrim}.
\begin{itemize}
\item For $(10,31)$ and $(11,35)$ all cubes are refuted with checked proofs. The hardest is
  $(11,35)$, $\Delta=7$, with $926$\,s to solve and $1645$\,s to check.
\item For $(12,41)$ the cubes $\Delta\in\{0,\dots,6,10,11\}$ are refuted with checked proofs.
  The cubes $\Delta=7,8,9$ remain, and we estimate them at a few hundred core-hours.
\end{itemize}

\begin{remark}
Two facts show how tight these thresholds are.
\begin{itemize}
\item At $(11,35)$ only $2$ of the $15{,}108{,}047$ graphs avoid a $12$-clique in $Q(G)$.
\item One of the $131$ extremal $11$-vertex graphs is already $10$-good.
\end{itemize}
Among the good $(12,40)$ graphs found, all have maximum degree $7$, $8$ or $9$. We found
$3{,}330$ with $\Delta\le9$, a count made by one method only.
\end{remark}

% =====================================================================
\section{Lower bounds for $13\le n\le128$}

Table~\ref{tab:lb} lists the best certificates found. The methods were tabu search on
(edge set, colouring) pairs; SAT with a prescribed semiregular automorphism acting on
vertices and colours; and the cyclic family of Section~\ref{sec:cyclic}. ``Type'' records
whether the certificate is rotation-equivariant on $\Z_n$ (cyclic), was found with an
imposed automorphism (symmetric), or neither (search). All entries are lower bounds only.

{\small
\begin{longtable}{rrrrl}
\caption{Lower bounds for $\M(n)$, with AlphaEvolve values from~\cite{TaoAE} for
  comparison.}\label{tab:lb}\\
\toprule
$n$ & AlphaEvolve & this note & $e/n^2$ & type\\
\midrule
\endfirsthead
\toprule
$n$ & AlphaEvolve & this note & $e/n^2$ & type\\
\midrule
\endhead
\bottomrule
\endlastfoot
% generated by make_tables.py -- do not edit
13 & 47 & 47 & 0.278 & search \\
14 & 54 & 55 & 0.281 & search \\
15 & 62 & 63 & 0.280 & symmetric \\
16 & 69 & 70 & 0.273 & symmetric \\
17 & 77 & 78 & 0.270 & symmetric \\
18 & 85 & 90 & 0.278 & cyclic \\
19 & 94 & 96 & 0.266 & symmetric \\
20 & 103 & 109 & 0.273 & search \\
21 & 112 & 126 & 0.286 & cyclic \\
22 & 123 & 133 & 0.275 & search \\
23 & 133 & 141 & 0.267 & search \\
24 & 143 & 150 & 0.260 & symmetric \\
25 & 155 & 160 & 0.256 & symmetric \\
26 & -- & 169 & 0.250 & symmetric \\
27 & -- & 189 & 0.259 & cyclic \\
28 & -- & 196 & 0.250 & cyclic \\
29 & -- & 210 & 0.250 & symmetric \\
30 & -- & 225 & 0.250 & symmetric \\
31 & -- & 248 & 0.258 & cyclic \\
32 & -- & 256 & 0.250 & cyclic \\
33 & -- & 264 & 0.242 & cyclic \\
34 & -- & 274 & 0.237 & search \\
35 & -- & 288 & 0.235 & search \\
36 & -- & 324 & 0.250 & cyclic \\
37 & -- & 333 & 0.243 & cyclic \\
38 & -- & 342 & 0.237 & cyclic \\
39 & -- & 351 & 0.231 & cyclic \\
40 & -- & 400 & 0.250 & cyclic \\
41 & -- & 410 & 0.244 & cyclic \\
42 & -- & 420 & 0.238 & cyclic \\
43 & -- & 430 & 0.233 & cyclic \\
44 & -- & 440 & 0.227 & cyclic \\
45 & -- & 495 & 0.244 & cyclic \\
46 & -- & 506 & 0.239 & cyclic \\
47 & -- & 518 & 0.234 & search \\
48 & -- & 528 & 0.229 & cyclic \\
49 & -- & 539 & 0.224 & cyclic \\
50 & -- & 550 & 0.220 & cyclic \\
51 & -- & 612 & 0.235 & cyclic \\
52 & -- & 624 & 0.231 & cyclic \\
53 & -- & 689 & 0.245 & cyclic \\
54 & -- & 690 & 0.237 & padded \\
55 & -- & 691 & 0.228 & padded \\
56 & -- & 728 & 0.232 & cyclic \\
57 & -- & 741 & 0.228 & cyclic \\
58 & -- & 754 & 0.224 & cyclic \\
59 & -- & 767 & 0.220 & cyclic \\
60 & -- & 780 & 0.217 & cyclic \\
61 & -- & 915 & 0.246 & cyclic \\
62 & -- & 916 & 0.238 & padded \\
63 & -- & 917 & 0.231 & padded \\
64 & -- & 918 & 0.224 & padded \\
65 & -- & 919 & 0.218 & padded \\
67 & -- & 938 & 0.209 & cyclic \\
68 & -- & 952 & 0.206 & cyclic \\
69 & -- & 966 & 0.203 & cyclic \\
70 & -- & 1050 & 0.214 & cyclic \\
71 & -- & 1065 & 0.211 & cyclic \\
72 & -- & 1080 & 0.208 & cyclic \\
73 & -- & 1095 & 0.205 & cyclic \\
74 & -- & 1110 & 0.203 & cyclic \\
75 & -- & 1125 & 0.200 & cyclic \\
76 & -- & 1216 & 0.211 & cyclic \\
78 & -- & 1248 & 0.205 & cyclic \\
79 & -- & 1264 & 0.203 & cyclic \\
80 & -- & 1280 & 0.200 & cyclic \\
81 & -- & 1296 & 0.198 & cyclic \\
82 & -- & 1312 & 0.195 & cyclic \\
83 & -- & 1411 & 0.205 & cyclic \\
84 & -- & 1428 & 0.202 & cyclic \\
85 & -- & 1445 & 0.200 & cyclic \\
86 & -- & 1462 & 0.198 & cyclic \\
88 & -- & 1496 & 0.193 & cyclic \\
89 & -- & 1519 & 0.192 & search \\
90 & -- & 1530 & 0.189 & cyclic \\
91 & -- & 1638 & 0.198 & cyclic \\
92 & -- & 1656 & 0.196 & cyclic \\
93 & -- & 1674 & 0.194 & cyclic \\
94 & -- & 1692 & 0.191 & cyclic \\
95 & -- & 1710 & 0.189 & cyclic \\
97 & -- & 1843 & 0.196 & cyclic \\
99 & -- & 1881 & 0.192 & cyclic \\
100 & -- & 1900 & 0.190 & cyclic \\
101 & -- & 1901 & 0.186 & padded \\
103 & -- & 1938 & 0.183 & search \\
106 & -- & 1941 & 0.173 & padded \\
107 & -- & 2065 & 0.180 & search \\
109 & -- & 2136 & 0.180 & search \\
113 & -- & 2277 & 0.178 & search \\
118 & -- & 2282 & 0.164 & padded \\
121 & -- & 2568 & 0.175 & search \\
122 & -- & 2569 & 0.173 & padded \\
125 & -- & 2716 & 0.174 & search \\
127 & -- & 2798 & 0.173 & search \\
128 & -- & 2831 & 0.173 & search \\

\end{longtable}
}

The ratio $e/n^2$ stays near $1/4$ up to $n\approx 40$ and then drifts down to about $0.17$
at $n=128$. Search effort per $n$ was much smaller at large $n$. Fitting decaying models
($cn^{2-\delta}$, $cn^2/\log n$, $cn^2e^{-a\sqrt{\log n}}$) to this data gives
extrapolations that disagree widely, so we draw no asymptotic conclusion from them.

% =====================================================================
\section{Cyclic orbit colourings}\label{sec:cyclic}

Let $H=\lceil n/2\rceil-1$. For $D\subseteq\{1,\dots,H\}$ and $g:D\to\Z_n$, define the
graph $\Gamma(n,D,g)$ on $\Z_n$:
\begin{itemize}
\item its edges are $\{i,i+d\}$ for $i\in\Z_n$ and $d\in D$;
\item the edge $\{i,i+d\}$ has colour $i+g(d)$.
\end{itemize}
It has $n|D|$ edges and uses $n$ colours, and rotation $i\mapsto i+1$ is an automorphism
that shifts every colour by $1$. Let $\kmax(n)$ be the largest $|D|$ for which some $g$
makes $\Gamma(n,D,g)$ $C_4$-rainbow. Then $\M(n)\ge n\,\kmax(n)$.

\begin{lemma}\label{lem:walk}
Write each $4$-cycle through $0$ as a closed walk $0\to v_1\to v_2\to v_3\to0$ on
distinct vertices, with signed steps $\pm d$, $d\in D$.
\begin{itemize}
\item A step $v\to v+d$ uses the edge $\{v,v+d\}$, of colour $v+g(d)$.
\item A step $v\to v-d$ uses $\{v-d,v\}$, of colour $(v-d)+g(d)$.
\end{itemize}
Thus each edge on the walk has colour $o+g(d)$ for an offset $o$ determined by the walk.
Two of its edges, with orbits $d_a\ne d_b$ and offsets $o_a,o_b$, have equal colours iff
$g(d_a)-g(d_b)\equiv o_b-o_a$. Edges in the same orbit never collide. Finally,
$\Gamma(n,D,g)$ is $C_4$-rainbow iff no such collision occurs on any $4$-cycle through $0$.
\end{lemma}

\begin{proof}
The colour formulas are the definition. Rotation by $t$ maps a $4$-cycle to a $4$-cycle and
adds $t$ to every colour, which preserves equal colours; this gives the reduction to cycles
through $0$. Two distinct edges in the same orbit have distinct offsets, so their colours
differ.
\end{proof}

\begin{lemma}\label{lem:mult}
Let $u\in\Z_n^\times$. For $d\in D$ write $ud\equiv\sigma e$ with $e\in\{1,\dots,H\}$ and
$\sigma=\pm1$. Define
\[
  g'(e)=\begin{cases} u\,g(d) & \text{if } \sigma=+1,\\ e+u\,g(d) & \text{if } \sigma=-1.\end{cases}
\]
If $\Gamma(n,D,g)$ is $C_4$-rainbow, then so is $\Gamma(n,D',g')$, where
$D'=\{e: d\in D\}$.
\end{lemma}

\begin{proof}
Apply the vertex map $i\mapsto ui$ and the colour map $c\mapsto uc$. Both are bijections, so
rainbow $4$-cycles go to rainbow $4$-cycles. The edge $\{i,i+d\}$ of colour $i+g(d)$ goes to
$\{ui,ui+ud\}$, of colour $ui+u\,g(d)$.
\begin{itemize}
\item If $ud\equiv e$, this is the edge $\{j,j+e\}$ with $j=ui$, of colour $j+u\,g(d)$.
\item If $ud\equiv -e$, it is $\{j,j+e\}$ with $j=ui-e$, of colour $j+e+u\,g(d)$.
\end{itemize}
\end{proof}

Together with the colour translations $g\mapsto g+t$, Lemma~\ref{lem:mult} gives the
symmetry group we use for symmetry breaking (lexicographic leaders of the orbit-presence
vector) and for listing solutions up to symmetry. We also checked the lemma numerically: all
$9{,}886$ images of the optimal solutions for $n\in\{20,21,24,27,28,31,32\}$ are valid.
Dropping the term $e$ in the case $\sigma=-1$ fails for every $u\ne1$ at $n=28$ and $40$.

\subsection{Exact optima}
Table~\ref{tab:cyc} gives $\kmax(n)$. Each exact value combines a solution at $\kmax(n)$
with an UNSAT result at $\kmax(n)+1$.
\begin{itemize}
\item For $n\le 36$ the UNSAT results are solver results, not DRAT-checked. For $n\le 28$ an
  independent backtracking counter (no SAT, no symmetry reduction) agrees with every count.
\item For $n=37,40,41$ the UNSAT proofs were checked by \texttt{drat-trim}.
\item For $n=38$ the LRAT proof was checked by a small checker written for the task. For
  $n=39$ the UNSAT result is unchecked.
\end{itemize}

{\small
\begin{longtable}{rrrl}
\caption{Cyclic orbit colourings: $\kmax(n)$ (exact) or the best $k$ found.}\label{tab:cyc}\\
\toprule
$n$ & $\kmax(n)$ & $k/n$ & status\\
\midrule
\endfirsthead
\toprule
$n$ & $\kmax(n)$ & $k/n$ & status\\
\midrule
\endhead
\bottomrule
\endlastfoot
% generated by make_tables.py -- do not edit
12 & 3 & 0.250 & exact (solver UNSAT; brute force) \\
13 & 3 & 0.231 & exact (DRAT, drat-trim) \\
14 & 3 & 0.214 & exact (DRAT, drat-trim) \\
15 & 4 & 0.267 & exact (DRAT, drat-trim) \\
16 & 4 & 0.250 & exact (DRAT, drat-trim) \\
17 & 4 & 0.235 & exact (DRAT, drat-trim) \\
18 & 5 & 0.278 & exact (DRAT, drat-trim) \\
19 & 5 & 0.263 & exact (DRAT, drat-trim) \\
20 & 5 & 0.250 & exact (DRAT, drat-trim) \\
21 & 6 & 0.286 & exact (DRAT, drat-trim) \\
22 & 5 & 0.227 & exact (DRAT, drat-trim) \\
23 & 6 & 0.261 & exact (DRAT, drat-trim) \\
24 & 6 & 0.250 & exact (DRAT, drat-trim) \\
25 & 6 & 0.240 & exact (DRAT, drat-trim) \\
26 & 6 & 0.231 & exact (DRAT, drat-trim) \\
27 & 7 & 0.259 & exact (DRAT, drat-trim) \\
28 & 7 & 0.250 & exact (DRAT, drat-trim) \\
29 & 7 & 0.241 & exact (DRAT, drat-trim) \\
30 & 7 & 0.233 & exact (DRAT, drat-trim) \\
31 & 8 & 0.258 & exact (DRAT, drat-trim) \\
32 & 8 & 0.250 & exact (DRAT, drat-trim) \\
33 & 8 & 0.242 & exact (DRAT, drat-trim) \\
34 & 8 & 0.235 & exact (DRAT, drat-trim) \\
35 & 8 & 0.229 & exact (DRAT, drat-trim) \\
36 & 9 & 0.250 & exact (DRAT, drat-trim) \\
37 & 9 & 0.243 & exact (DRAT, drat-trim) \\
38 & 9 & 0.237 & exact (DRAT, drat-trim) \\
39 & 9 & 0.231 & exact (DRAT, drat-trim) \\
40 & 10 & 0.250 & exact (DRAT, drat-trim) \\
41 & 10 & 0.244 & exact (DRAT, drat-trim) \\
42 & 10 & 0.238 & exact (two solvers, proof unchecked) \\
43 & $\ge 10$ & 0.233 & lower bound \\
44 & 10 & 0.227 & exact (LRAT, \texttt{cake\_lpr}, all leaves) \\
45 & $\ge 11$ & 0.244 & lower bound \\
46 & $\ge 11$ & 0.239 & lower bound \\
47 & 11 & 0.234 & exact (cube-and-conquer, proofs unchecked) \\
48 & $\ge 11$ & 0.229 & lower bound \\
49 & $\ge 11$ & 0.224 & lower bound \\
50 & $\ge 11$ & 0.220 & lower bound \\
51 & $\ge 12$ & 0.235 & lower bound \\
52 & $\ge 12$ & 0.231 & lower bound \\
53 & $\ge 13$ & 0.245 & lower bound \\
54 & $\ge 12$ & 0.222 & lower bound \\
55 & $\ge 12$ & 0.218 & lower bound \\
56 & $\ge 13$ & 0.232 & lower bound \\
57 & $\ge 13$ & 0.228 & lower bound \\
58 & $\ge 13$ & 0.224 & lower bound \\
59 & $\ge 13$ & 0.220 & lower bound \\
60 & $\ge 13$ & 0.217 & lower bound \\
61 & $\ge 15$ & 0.246 & lower bound \\
62 & $\ge 13$ & 0.210 & lower bound \\
63 & $\ge 13$ & 0.206 & lower bound \\
65 & $\ge 14$ & 0.215 & lower bound \\
67 & $\ge 14$ & 0.209 & lower bound \\
68 & $\ge 14$ & 0.206 & lower bound \\
69 & $\ge 14$ & 0.203 & lower bound \\
70 & $\ge 15$ & 0.214 & lower bound \\
71 & $\ge 15$ & 0.211 & lower bound \\
72 & $\ge 15$ & 0.208 & lower bound \\
73 & $\ge 15$ & 0.205 & lower bound \\
74 & $\ge 15$ & 0.203 & lower bound \\
75 & $\ge 15$ & 0.200 & lower bound \\
76 & $\ge 16$ & 0.211 & lower bound \\
78 & $\ge 16$ & 0.205 & lower bound \\
79 & $\ge 16$ & 0.203 & lower bound \\
80 & $\ge 16$ & 0.200 & lower bound \\
81 & $\ge 16$ & 0.198 & lower bound \\
82 & $\ge 16$ & 0.195 & lower bound \\
83 & $\ge 17$ & 0.205 & lower bound \\
84 & $\ge 17$ & 0.202 & lower bound \\
85 & $\ge 17$ & 0.200 & lower bound \\
86 & $\ge 17$ & 0.198 & lower bound \\
88 & $\ge 17$ & 0.193 & lower bound \\
89 & $\ge 17$ & 0.191 & lower bound \\
90 & $\ge 17$ & 0.189 & lower bound \\
91 & $\ge 18$ & 0.198 & lower bound \\
92 & $\ge 18$ & 0.196 & lower bound \\
93 & $\ge 18$ & 0.194 & lower bound \\
94 & $\ge 18$ & 0.191 & lower bound \\
95 & $\ge 18$ & 0.189 & lower bound \\
97 & $\ge 19$ & 0.196 & lower bound \\
99 & $\ge 19$ & 0.192 & lower bound \\
100 & $\ge 19$ & 0.190 & lower bound \\

\end{longtable}
}

\begin{remark}[$n\equiv0\pmod 4$]
For $n=12,16,20,24,28,32,36,40$ we have exactly $\kmax(n)=n/4$, and hence
$\M(n)\ge n^2/4$. The number of solutions at $k=n/4$, up to the group generated by
multipliers and colour translations, is
\[
  24,\ 58,\ 165,\ 362,\ 285,\ 128,\ 59,\ 4 \qquad (n=12,16,\dots,40).
\]
These counts come from complete SAT enumeration with blocking clauses. They peak at $n=24$
and fall sharply after that. No solution with $k=n/4$ has been found for any multiple of $4$
in $[44,100]$.
% TODO(final CYC wave, 2026-10-25 10:15Z): update with k_max(44) if decided.
For several odd $n$, $\kmax(n)/n$ exceeds $1/4$ (for example $\kmax(21)=6$ and $\kmax(31)=8$).
For $n=18,21,31$ the optimal solution is unique up to symmetry.
\end{remark}

\subsection{Structured rules give $o(n)$ orbits}

The optimal solutions found show no low-degree polynomial structure. Natural rules for $g$
also fail. The following proposition explains why for linear rules.

\begin{proposition}\label{prop:linear}
Fix integers $u,v$ with $v\ge1$ and $b=u/v\notin\{0,1\}$. Let $n$ be coprime to $v$.
Suppose $\Gamma(n,D,g)$ is $C_4$-rainbow and $g(d)\equiv bd\pmod n$ for all $d\in D$.
Then $D$ contains no set $\{m,\ m+vr,\ m+(v-u)r,\ m+ur\}$ with $r\ge1$,
$ur\not\equiv0$ and $(v-u)r\not\equiv0\pmod n$.
Consequently $|D|=o(n)$.
\end{proposition}

\begin{proof}
Let $d_1=m+vr$, $d_2=m+(v-u)r$, $d_3=m$ and $d_4=m+ur$, all in $D$. Then
$d_1-d_2+d_3-d_4=0$.
\begin{itemize}
\item The walk with steps $+d_1,-d_2,+d_3,-d_4$ visits $0$, $d_1$, $d_1-d_2=ur$ and
  $ur+m=d_4$.
\item These four vertices are distinct. First, $d_1,d_4\not\equiv0$, and $ur\not\equiv0$.
  Next, $d_1-d_4=(v-u)r\not\equiv0$ and $d_1-ur=d_2\not\equiv0$. Finally,
  $d_4-ur=m\not\equiv0$.
\item The first edge has colour $g(d_1)=b(m+vr)=bm+ur$, using $bv\equiv u$. The third edge
  leaves $ur$ with step $+d_3$, so its colour is $ur+g(d_3)=ur+bm$.
\end{itemize}
The two colours are equal, which contradicts rainbowness.

For the consequence, note that $\{0,v,v-u,u\}$ is a fixed finite pattern. By Szemer\'edi's
theorem~\cite{Sz75}, every subset of $\{1,\dots,H\}$ of density $\delta>0$ contains a copy
$m+r\cdot\{0,v,v-u,u\}$ with $r\ge1$. By Varnavides' averaging~\cite{Var59} it contains
$\gg_\delta n^2$ such copies. Only $O(n)$ of them violate the non-degeneracy conditions:
at most $|u|+|u-v|$ residues $r$ are excluded, each with at most $n$ choices of $m$.
\end{proof}

\begin{remark}
For $b=2$ the forbidden pattern is a $4$-term arithmetic progression, so
$|D|\le r_4(H)$; by the bound of Green and Tao~\cite{GT17} this is
$O(n(\log n)^{-c})$.
The proof uses only $g(d_1)$ and $g(d_3)$. Hence the same conclusion holds for any
$A\subseteq D$ on which $g(d)\equiv bd+c$ for a fixed $c$. So rules in which $g$ is piecewise
affine on $D$, with boundedly many pieces and slopes of bounded height, give $|D|=o(n)$.
For slope $b\in\{0,1\}$ such a piece has at most one element.
\end{remark}

\begin{remark}[Experiments]
Two further experiments point the same way. Both are evidence, not theorems.
\begin{itemize}
\item \emph{Coset rules.} Take $D=\pm cK$ for an odd-order subgroup $K\le\Z_p^\times$; then for
  colourings invariant under $K$ the colour of an edge is an affine function of its endpoints
  (for a suitable orientation). These rules achieve
  $k=(p-1)/4$ at $p=13,29,37,53,61$ and give the new bounds $\M(53)\ge689$ and
  $\M(61)\ge915$. Every fixed index eventually fails, and the best such $k$ grows roughly
  like $p^{3/4}$.
\item \emph{Doubling.} Lifting an optimal $D$ at $m=12,16,20$ to
  $D\cup(m-D)\subseteq\Z_{2m}$ never admits a valid $g$.
\end{itemize}
\end{remark}

% =====================================================================
\section{Other families}

Sum colourings on $\Z_n$ give at most $|E|\approx r_3(n)\,n/2=o(n^2)$ edges, by Roth's theorem:
$\{i,j\}$ is an edge iff $i+j\in S$, with colour $i+j$. In our range they achieve only
$0.10$--$0.13\,n^2$. Allowing an arbitrary edge set with colour $i+j$ gives about
$0.17$--$0.21\,n^2$. Bipartite affine colourings over $\mathbb F_p$ use $p-1$ colours on
$2p$ vertices, the colour class of slope $a$ being a partial map $x\mapsto ax+\beta(a)$. The
best we found have $m/p^2$ falling from about $0.48$ at $p=5$ to about $0.30$ at $p=101$,
that is about $0.075\,n^2$. These relate to the restricted results in~\cite{Repo}.
The construction of~\cite{BEGS89} from $(7,4)$-free triple systems was also weaker in our
range.

% =====================================================================
\section{Open problems}

\begin{enumerate}
\item Determine $\M(13)$; we know $\M(13)\ge 47$.
  \begin{itemize}
  \item Pruned enumeration of the $584{,}089{,}835{,}857$ classes of $(13,48)$ graphs is
    estimated at $10^4$ CPU-hours.
  \item Splitting by maximum degree and using SAT cubes for large $\Delta$ may be cheaper.
  \end{itemize}
\item Give a fully certified proof of $\M(12)=40$, by checking the cubes $\Delta=7,8,9$.
\item Is $\kmax(n)=o(n)$? A construction with $\kmax(n)\ge\varepsilon n$ would answer
  Erd\H{o}s \#810 positively. Is $\kmax(4m)<m$ for some $m$? The first open case is $n=44$.
\end{enumerate}

% =====================================================================
\appendix
\section{Reproducibility}\label{app:repro}

The repository~\cite{Repo}, branch \texttt{bruteforce-2026-10}, contains:
\begin{itemize}
\item every certificate, as JSON lists of $(u,v,\text{colour})$ or, for cyclic colourings,
  of $(d,g(d))$ pairs;
\item the enumeration and SAT programs;
\item solver logs and the hashes of the CNF files;
\item proof-checker outputs.
\end{itemize}
The script \texttt{paper/make\_tables.py} re-checks every certificate with a short checker
that enumerates all $4$-cycles, and regenerates Tables~\ref{tab:lb} and~\ref{tab:cyc}.
% TODO(author): pin final commit hashes here.

% TODO(author): acknowledgements; disclosure of AI assistance in the computations and in
% drafting, following the venue's policy.

\bibliographystyle{amsplain}
\bibliography{refs}

\end{document}
